\begin{center}
\textbf{RESUMO}
\end{center}

%
% O resumo deve ser organizado em apenas um parágrafo mesmo.
% O número de folha é o número de páginas do PDF -2. Isto ocorre pois na versão final (capa dura) a capa é removida e as duas primeiras páginas são impressas em uma % folha apenas (frente e verso).
%

$\!$\\

\hspace{-1.3cm}\textbf{RABELO DE TOLEDO}, Igor. \textit{Otimização de Portfólios Financeiros Utilizando Inteligência Computacional}. XX f. Trabalho de Conclusão de Curso (Bacharel em Engenharia Elétrica com ênfase em Sistemas e Computação) - Faculdade de Engenharia, Universidade do Estado do Rio de Janeiro~(UERJ), Rio de Janeiro, 2026.

\vspace{.2cm}

Métodos tradicionais de alocação de investimentos frequentemente geram carteiras instáveis por dependerem apenas de médias históricas. Essa limitação reforça a importância de utilizar abordagens que tornem o processo de alocação mais consistente diante das incertezas do mercado financeiro. Nesse contexto, este trabalho propõe um sistema automatizado para otimizar um portfólio de até 20 ações de alta liquidez da bolsa brasileira (B3), utilizando o modelo de Black-Litterman,  que combina previsões matemáticas de retorno geradas por aprendizado de máquina (XGBoost) e a análise do sentimento de manchetes financeiras feita por inteligência artificial (GPT), durante a construção de projeções de investimento.
Essas projeções são integradas ao equilíbrio natural do mercado de forma inteligente: o peso de cada decisão é ajustado por um controle dinâmico de confiança, que considera o histórico de acertos do modelo e o volume de notícias disponíveis no dia. O sistema opera de ponta a ponta — desde a coleta de dados até a alocação final —, permitindo apenas posições de compra e impondo um limite máximo de concentração de capital por empresa.
Avaliado em simulações com dados não vistos durante o treinamento e considerando custos reais de transação, o modelo foi comparado a abordagens clássicas (Markowitz) e à taxa Selic. Os resultados indicam que integrar inteligência artificial, leitura automatizada de textos e controle estatístico de risco é uma via prática e promissora para apoiar a tomada de decisão na B3.

\vspace{1cm}

\hspace{-1.3cm}Palavras-chave: Otimização de Portfólio; Modelo de Black-Litterman; Aprendizado de Máquina; Análise de Sentimento; Mercado de Ações.